\section{Normalized transitions on the parabolic tails}
\label{sec:source-normalization}

Above the coefficient height scale, the infinity transition consists
of two parabolic tails joined by a bounded connector. Normalization by
each tail's exact value and first variation gives convergent nonlinear
maps, while the common singular power cancels in the ratio of the two
absolute variations. This gives the signed high-height estimate of
Corollary~\ref{rew:cor-high-derivative-signed} for the actual return
\eqref{eq:intro-displacement}.

All constants are uniform on the stated compact parameter sets.
A positive unit is bounded above and away from zero;
\(O_{C^k}\) denotes ordinary phase derivatives through order \(k\)
unless weights are displayed. Choose the finite intrinsic normal-form
regularity before constructing the maps. Exterior observers use
ordinary \(C^3\) bounds on their actual buffered images.

Put \(\alpha _0=s^2a,\ \alpha _1=sp,\ v=\gamma+sX,\ z=s^2Y\).
For \(s>0\), a positive change of time gives
\begin{equation}\label{ann:source-field}
\begin{split}
 X'&=a+pX-Y+A(\gamma+sX)Y-3\gamma X^2-sX^3
                 -C(\gamma+sX)^2Y,\\
 Y'&=-Y\{(1+2\gamma)X+sX^2+Cs(\gamma+sX)Y\}.
\end{split}
\end{equation}
Here \(s\) is a general positive scale. For negative opening one may
take \(s=\sqrt{-\alpha _0}\), whereas normalizing the height of a
reference orbit allows \(a\) to have either sign.

\subsection{The limiting transition}
\label{ann:subsec-source-scattering}

At \(s=\gamma=0\), the equation is
\begin{equation}\label{ann:limiting-source}
 X'=a+pX-Y,\qquad Y'=-XY.
\end{equation}
We call a complete orbit \emph{strict} if it is a graph on \(\R\)
with \(Y>0\) and \(D=Y-a-pX>0\). Positive time then corresponds to
decreasing \(X\). A compact strict family means a compact family of
such graphs, with uniform positive margins for these two inequalities
on each fixed compact \(X\)-interval.

Removing the common quadratic, linear and logarithmic terms at each end
gives finite labels for these complete graphs.

\begin{lemma}\label{ann:lem-renormalized-ends}
On a compact strict family the limits
\begin{equation}\label{eq:renormalized-ends}
 B_\pm=\lim_{X\to\pm\infty}
 \{Y-X^2/2-2pX-(2a-4p^2)\log|X|\}
\end{equation}
exist. They give increasing coordinates on the family and depend
real-analytically on the parameters and the orbit. The corresponding
maps and their inverses have bounded derivatives of every prescribed
finite order on smaller compact strict families.
\end{lemma}
\begin{proof}
Bound \(|a|,|p|\) by \(M\ge1\). On either outward branch the graph
equation \(Y_X=XY/(Y-a-pX)\), with \(r=|X|\), gives
\(Y_r\ge rY/(Y+M(1+r))\). Below \(Y=M(1+r)\),
\((\log Y)_r\ge1/(4M)\) for \(r\ge1\); above \(r=2M\)
this line is a lower barrier. Compactness supplies a common entrance;
after enlarging the cut, \(Y\ge r^2/8\) and \(Y_r=r+O(1)\), uniformly.
Writing \(c=2a-4p^2\) and
\(Y=X^2/2+2pX+c\log|X|+\mathcal B\), direct expansion in a
parabolic corridor gives
\[
 |\mathcal B_X|\le
 C\frac{1+\log|X|+|\mathcal B|}{X^2},
 \qquad |\partial_{\mathcal B}\mathcal B_X|\le C X^{-2}.
\]
Gronwall's inequality gives boundedness and convergence. Beyond a fixed
large cut, the prescribed-limit integral equation is a contraction on a
common complex neighborhood, since its denominator is comparable to
\(X^2\). Holomorphic Picard iteration and finite-interval dependence
give analyticity. The positive scalar first variation and its reciprocal
are bounded, proving the coordinate and inverse assertions.
\end{proof}

Denote the scattering map by
\begin{equation}\label{ann:limiting-scattering-map}
 S_{a,p}(b)=B_-,\qquad b=B_+.
\end{equation}
Its curvature has a particularly simple sign.

\begin{proposition}\label{ann:prop-source-curvature}
On every strict component,
\begin{equation}\label{ann:source-scattering-division}
 S_{a,p}(b)=b-p\mathcal I(a,p,b),\qquad
 \mathcal I_{bb}>0.
\end{equation}
In particular, the curvature of \(S_{a,p}\) has sign
\(-\sgn p\) when \(p\ne0\).
\end{proposition}
\begin{proof}
When \(p=0\), reflection \(X\mapsto-X\) reverses the flow and gives
\(S_{a,0}=\mathrm{id}\). Analytic division in \(p\) defines \(\mathcal I\).
For \(\sigma=\partial_bY>0\) and \(\tau=\partial_b^2Y\), the graph
equation gives
\[
 \sigma_X=-\frac{X(a+pX)}{D^2}\sigma,\qquad
 \tau_X=-\frac{X(a+pX)}{D^2}\tau
              +\frac{2X(a+pX)}{D^3}\sigma^2.
\]
Integrating the equation for \(\tau/\sigma\) from the incoming to the
outgoing end yields
\[
 \frac{S''}{S'}=-2\int_\R\frac{X(a+pX)\sigma}{D^3}\,dX.
\]
To evaluate this integral, use
\[
 D_X=X-p+\frac{X(a+pX)}D,\qquad
 (\sigma/Y)_X=-X\sigma/D^2.
\]
The integral of \((\sigma/D)_X\), followed by that of
\((\sigma/Y)_X\), therefore gives
\begin{equation}\label{ann:source-log-curvature}
 \frac{S''(b)}{S'(b)}
       =-p\int_\R\frac{\sigma(X,b)}{D(X,b)^2}\,dX.
\end{equation}
All boundary terms vanish by \eqref{eq:renormalized-ends}.
Thus
\(\mathcal I_{bb}=S'\int_\R\sigma D^{-2}\,dX>0\), also at \(p=0\)
by continuity.
\end{proof}

\subsection{Normalization at a reference orbit}

Fix a strict reference orbit and label nearby orbits by \(b\), with
reference value \(b_0\). For positive \(s\), let \(h(b)\) be their
height on \(\Sigma_+\), and put
\[
 k(b)=\mathcal C_{\theta,m}\bigl(T_\theta(h(b))\bigr).
\]
Thus \(h-a_\theta-mk=0\) is exactly the periodic-orbit equation.
The maps \(h,k\) include all transitions to the fixed sections. Their
slopes may vanish at different rates, so normalize their values and first
derivatives separately:
\begin{equation}\label{eq:reference-normalization}
 H(b)=b_0+\frac{h(b)-h(b_0)}{h'(b_0)},\quad
 K(b)=b_0+\frac{k(b)-k(b_0)}{k'(b_0)},\quad
 \Phi=K\circ H^{-1}.
\end{equation}
The input \(\beta=H(b)\) is affine in the actual incoming height.
Reversing the sign of the return equation and dividing by \(mk'(b_0)>0\)
gives exactly
\begin{equation}\label{eq:affine-normalized-displacement}
 D(\beta)=c_0+c_1\beta+\Phi(\beta)=0 .
\end{equation}
The constants retain both centers and both slopes. No bound on their
ratio is assumed.

\begin{proposition}[Convergence on the actual sections]
\label{prop:source-normalization}
\label{ann:prop-source-curvature-persistence}
On each compact strict family, \(\Phi\) extends continuously in
phase \(C^3\) to \(s=\gamma=0\), with limit the positive affine
normalization of \(S_{a,p}\) in \eqref{eq:reference-normalization}.
Let \(\Phi^{\rm int}\) denote the normalized transition before the
exterior observers. Its maps and inverses have every prescribed finite
phase order; near \(p=0\), when the division intervals remain strict,
its \(s\)-derivative for \(s\ge0\) and \(p\)-derivative at \(s=0\)
extend in each prescribed \(C^k\), using one extra intrinsic phase
order. Moreover, on common enlarged domains,
\[
 \|\Phi-\Phi^{\rm int}\|_{C^3}=O(s^{2-\nu}),\qquad 0<\nu<1.
\]
Consequently the actual \(s\)-secant extends continuously in
\(C^3\) at zero. At \(p=s=0\), for small \(\gamma\),
\(\Phi=\mathrm{id}\). All assertions are uniform in passive
coefficients and \(0\le m\le m_0\).
\end{proposition}

\begin{proof}
Ordinary ODE dependence controls the compact core. On each outward tail,
set \(X=\sigma r\), \(\sigma=\pm1\); then
\begin{equation}\label{eq:parabolic-tail}
 Y_r=F_\sigma(r,Y)=
 \frac{(1+2\gamma+\sigma sr)rY+
       \sigma Cs(\gamma+\sigma sr)Y^2}
 {\ell_\sigma Y+3\gamma r^2+\sigma sr^3-\sigma pr-a},
 \quad
 \ell_\sigma=1-A(\gamma+\sigma sr)+C(\gamma+\sigma sr)^2 .
\end{equation}
Choose a small raw radius \(\delta\), then a large core radius \(L\),
then the bounds on \(s,\gamma\). The tail ends at \(R_s=\delta/s\),
where uniformly regular maps join \(v=\gamma\pm\delta\) to \(\Sigma_\pm\).

Write
\[
 L_\gamma=1-A\gamma+C\gamma^2,\quad
 W_\gamma=\frac{1-4\gamma}{2L_\gamma},\quad
 \lambda_\gamma=\frac{12\gamma}{1+2\gamma}.
\]
At \(s=p=a=0\), \(W=Y/r^2\) has equilibrium \(W_\gamma\)
and radial exponent \(-2+\lambda_\gamma\).
Fix \(0<\nu<1\), with \(|\lambda_\gamma|<\nu\). A small band about
\(W_\gamma\) is invariant; its logarithmic equation has perturbation
\(O(sr+|p|/r+|a|/r^2)\). Variation of constants gives
\begin{equation}\label{eq:parabolic-corridor}
 W-W_\gamma=O(sr+|p|/r+r^{-2+\nu}),\qquad L\le r\le R_s .
\end{equation}
This includes the logarithmic term for bounded \(a\) and makes the
denominator comparable to \(r^2\), uniformly on the compact strict family
after choosing \(L\).

Let \(Y_*(r)=Y(r,b_0)\) and \(j(r)=Y_b(r,b_0)>0\).
For the rational expression \(F=(\alpha Y^2+\beta Y)/(\ell Y+m)\),
polynomial division gives the exact identity
\begin{equation}\label{tail:rational-division}
 F_{Y^h}=(-1)^h h!\,
 \frac{m(\alpha m-\beta\ell)\ell^{h-2}}
      {(\ell Y+m)^{h+1}},\qquad h\ge2.
\end{equation}
Using the coefficients in \eqref{eq:parabolic-tail} and the corridor gives
\begin{align}
 F_Y-\lambda_\gamma/r
   &=O(s+|p|r^{-2}+r^{-3+\nu}),\label{ann:tail-first-derivative}\\
 |F_{Y^j}|&\le C_j\{sr^{-2j+2}+|\gamma|r^{-2j+1}
                    +|p|r^{-2j}+|a|r^{-2j-1}\},\quad j\ge2.
 \label{ann:tail-higher-derivatives}
\end{align}
At fixed \(r,Y\), their explicit coefficient derivatives satisfy
\begin{equation}\label{tail:rational-parameter-bounds}
 |\partial_s^{\rm exp}F_{Y^h}|\le C_hr^{-2h+2},\qquad
 |\partial_p^{\rm exp}F_{Y^h}|_{s=0}\le C_hr^{-2h+1},\quad h\ge2.
\end{equation}
Integrating \eqref{ann:tail-first-derivative}, using \(sR_s=\delta\), gives
\begin{equation}\label{ann:tail-first-variation}
 c r^{\lambda_\gamma}\le j(r)\le C r^{\lambda_\gamma}.
\end{equation}
The parameter variations, at fixed \(r\), satisfy
\begin{equation}\label{eq:tail-parameter-variations}
 |Y_{*,s}|\le Cr^3,\quad |(\log j)_s|\le Cr,\qquad
 |Y_{*,p}|_{s=0}\le Cr,\quad |(\log j)_p|_{s=0}\le C .
\end{equation}
For the height rows, apply variation of constants to
\(F_s=O(r^2)\), \(F_p|_{s=0}=O(1)\), with propagator
\eqref{ann:tail-first-variation}. For the slope rows,
\((\log j)_{\ell r}=F_{Y\ell}+F_{YY}Y_{*,\ell}\) is respectively
\(O(1)\) and \(O(r^{-2})\).

Subtracting the exact affine part gives
\begin{equation}\label{eq:normalized-difference}
 \eta=(Y-Y_*)/j,\qquad \eta_r=\mathcal N(r,\eta),
\end{equation}
with kernel \eqref{tail:normalized-kernel}. On a bounded
\(\eta\)-interval this remainder can also be written exactly as
\begin{equation}\label{tail:exact-affine-remainder}
 \mathcal N=
 \frac{j\eta^2 m(\alpha m-\beta\ell)}
 {D_*^2(D_*+\ell j\eta)},\qquad D_*=\ell Y_*+m.
\end{equation}
Both denominators are comparable to \(r^2\). The phase and parameter
formula \eqref{deg:det-parameter-phase-derivative}, together with
\eqref{ann:tail-higher-derivatives},
\eqref{tail:rational-parameter-bounds} and
\eqref{eq:tail-parameter-variations}, gives
\begin{align}
 \|\mathcal N\|_{C^k}
 &\le C\{sr^{-2+\nu}+|\gamma|r^{-3+\nu}
                       +|p|r^{-4+\nu}+|a|r^{-5+\nu}\},\label{eq:N-bound}\\
 \|\mathcal N_s\|_{C^k}&\le Cr^{-2+\nu},\qquad
 \|\mathcal N_p\|_{C^k,s=0}\le Cr^{-3+\nu}.
 \label{eq:N-parameter-bound}
\end{align}
The reference and slope are chosen by finite ODE dependence at the fixed
core cut. Thus \(Y_{*,\ell}\), \(j_\ell\) include their motion while
the normalized initial phase stays fixed. Every further phase row adds
\(j/r^2=O(r^{-2+\nu})\), so the bounds include the extra order needed
for the \(C^k\)-valued parameter variation.

These are the integrable majorants of Lemma~\ref{lem:scalar-normalization}.
The moving endpoint contributes
\[
 \|\mathcal N(R_s,\cdot)R_s'\|_{C^k}=O(s^{1-\nu})\longrightarrow0 .
\]
Finite-interval ODE dependence supplies local kernel convergence.
The lemma gives joint convergence, common domains and inverse bounds.

For each exterior map \(R\), its actual interval is
\(z=z_*+d_s\eta\), where \(d_s=s^2j(R_s)=O(s^{2-\nu})\).
Lemma~\ref{phys:lem-c3-adapters} gives on the admitted buffered image
\begin{equation}\label{eq:exterior-normalization}
 \frac{R(z_*+d_s\eta)-R(z_*)}{R'(z_*)d_s}
                  =\eta+O_{C^3}(d_s),
\end{equation}
including inverse and divided closing. With these observers denoted by
\(E_\pm\), the exact composition is
\(\Phi=E_-\circ\Phi^{\rm int}\circ E_+^{-1}\). The intrinsic
\(C^4\) bound and Lemma~\ref{phys:lem-c3-core-transport} give the
claimed \(O_{C^3}(s^{2-\nu})\) error. Its quotient by \(s\) tends to
zero, requiring only phase derivatives of the exterior maps.

At \(p=s=0\) a first integral is
\begin{equation}\label{ann:gamma-first-integral}
\begin{split}
 B_{a,\gamma}(X,Y)
 &=-Y^{-\nu_\gamma}
       \left(\frac{X^2}{2}-\frac{L_\gamma}{1-4\gamma}Y\right)
       -\frac{a}{1+2\gamma}\log Y\,
          \frac{1-e^{-\nu_\gamma\log Y}}{\nu_\gamma\log Y}
       -a\log2,\\
 &\hspace{35mm}\nu_\gamma=\frac{6\gamma}{1+2\gamma},
\end{split}
\end{equation}
with the quotient continued at zero. Its \(Y\)-derivative is positive
on the strict core. Use it as the initial orbit label, shifted to \(b_0\)
at the reference orbit. The same label occurs on the two tails.
Implicit differentiation gives \(Y_{bb}/Y_b=O(r^{-2+\nu})\),
and the analogous higher-order estimates. Therefore both normalized
tail limits are the identity. Together with
\eqref{eq:exterior-normalization} this proves \(\Phi=\mathrm{id}\).
At \(\gamma=0\), the limiting tail labels for general \(p\) are
\eqref{eq:renormalized-ends}, which proves the stated scattering limit.
\end{proof}

During intrinsic parameter integration, hold \(m\) independently
fixed in \([0,m_0]\). Its divided exterior map extends to both signs
of the finite opening; only the actual \(m(\theta)>0\) enters the
cycle equation.

\subsection{Exact division in the perturbation parameters}

On common enlarged strict domains, put
\begin{equation}\label{ann:exact-source-division}
 \begin{gathered}
 \Phi(\beta)=\beta+pU(\beta)+sV(\beta),\qquad
 U=\int_0^1\Phi_p^{\rm int}(a,tp,0,\gamma;\beta)\,dt,\\
 V=\int_0^1\Phi_s^{\rm int}(a,p,ts,\gamma;\beta)\,dt
          +\frac{\Phi-\Phi^{\rm int}}s .
 \end{gathered}
\end{equation}
The two integrals divide successively in \(p\) and \(s\), retaining
the actual exterior remainder. Both secants have continuous bounded
\(C^3\) limits: the first uses only the intrinsic
\(p\)-derivative at \(s=0\); the last quotient is
\(O_{C^3}(s^{1-\nu})\). The strict margins hold throughout both
intrinsic integration intervals, with fixed references and common
inverse images.

\subsection{Cancellation of the absolute tail derivatives}
\label{ann:subsec-high-source}

Use the actual affine incoming coordinate of Lemma~\ref{lem:closing}
and write
\[
 \Pi(x)=\mathcal C_{\theta,m}\bigl(T_\theta(a_\theta+x)\bigr),
 \qquad \Delta(x)=x-m\Pi(x).
\]
The fixed raw ports, their regular bridges, and the divided closing map
are all included in this same transition.

\begin{proposition}\label{sharp:upper-high}
On any fixed compact strict source family with its two outward tails,
the complete transition satisfies \(0<\Pi'(x)\le C\), uniformly for
small positive \(s\), small \(\gamma\), and \(0\le m\le m_0\).
The constant may depend
on the compact family. This includes a small signed \(a\)-block about
\(a=p=0\), with central height near one.

On the negative-opening source family \(a=-1\), there are constants
\(K_H,z_*,C_H>0\) such that, for either sign of small \(\gamma\),
\[
 Y_0>K_H(1+p^2),\qquad 0<s^2Y_0<z_*
 \quad\Longrightarrow\quad 0<\Pi'(x)\le C_H.
\]
Here \(C_H\) is independent of \(p\) and of the tested phase point.
After one common multiplier shrink, \(\Delta'>1/2\) on this whole
high-height interval.
\end{proposition}
\begin{proof}
Fix a compact strict family and a core cut \(L_t\). With
\(y_c=Y(0)\), ordinary first variation on \([-L_t,L_t]\) gives
\[
 A_\sigma=\frac{\partial Y(\sigma L_t;y_c)}{\partial y_c}
 \in[c_A,C_A],\qquad \sigma=\pm.
\]
The coefficient scale \(s\) is held fixed in this differentiation.
Equations \eqref{ann:tail-first-derivative} and
\eqref{ann:tail-first-variation}, before affine normalization, give on
each outward tail
\[
 J_\sigma(R)=(R/L_t)^{\lambda_\gamma}U_\sigma(R),\qquad
 \lambda_\gamma=\frac{12\gamma}{1+2\gamma},\qquad
 0<c\le U_\sigma,U_\sigma^{-1}\le C.
\]
At ports \(v=\gamma\pm d_H\), \(R=d_H/s\), the raw variations are
\begin{equation}\label{sharp:upper-tail-ratio}
 \frac{dz_\sigma}{dy_c}
 =s^{2-\lambda_\gamma}(d_H/L_t)^{\lambda_\gamma}
                                      A_\sigma U_\sigma(R),\qquad
 T'(z_+)=\frac{A_-U_-(R)}{A_+U_+(R)}.
\end{equation}
Here \(T\) is the passage between the raw ports. The common power
cancels; distinct fixed radii add only the bounded factor
\((d_-/d_+)^{\lambda_\gamma}\). Regular bridges are derivative units, and
\(\mathcal C_{\theta,m}'=f_\theta'(md_\theta)d_\theta'\) is a
uniform positive unit by Lemma~\ref{lem:closing}. Thus the derivative
of the complete transition in the actual incoming coordinate is
bounded on the fixed compact strict family.

For the signed block \(|a|,|p|\le1/8\),
\(3/4\le y_c\le5/4\), the half-graphs satisfy
\[
 3/4+r^2/4\le Y(r)\le5/4+r^2,\qquad
 Y-a-\sigma pr\ge39/64.
\]
These follow by scalar comparison; below \(Y/r^2=1/4\) the
logarithmic vector points upward, and above \(3/4\) it points
downward beyond a common fixed cut. Thus this is a compact strict
family with uniform parabolic entrances.

For the negative-opening high-height range, freeze
\(q=Y_0^{-1/2}\) at a selected point and rescale the source field to
\[
 (\bar a,\bar p,\bar\gamma,\bar s)
   =(-Y_0^{-1},pY_0^{-1/2},\gamma,sY_0^{1/2}),\qquad y(0)=1.
\]
Large \(K_H\), followed by small \(z_*\) and a coefficient shrink,
places these parameters in one compact strict source block about
\(y=1+x^2/2\). Apply the preceding bound with scale \(\bar s\).
The same constant \(C_H\) then works at every tested point, and
\(mC_H<1/2\) gives \(\Delta'>1/2\).
The scale is frozen during the phase derivative, with no bound on
\(\gamma\log\bar s\).
\end{proof}

\subsection{Extension through the height cap}

The derivative bound joins the low-height count to the fixed raw cap.
The overlap condition in the next lemma ensures that this continuation
introduces no additional critical point.

\begin{lemma}[Extension through the height cap]
\label{sharp:upper-height-splice}
On one ordered physical carrier, suppose a connected lower interval
has a prescribed critical-point bound and overlaps a bounded-height
interval where the complete transition derivative is uniformly bounded.
Write its scaled height as \(\eta=\tau Y_0\), \(\tau>0\), and
suppose these two intervals cover \(0<\eta\le\eta_{\max}\).
If Proposition~\ref{sharp:upper-high} applies above
\(Y_0=K_H(1+p^2)\) and
\begin{equation}
 \tau K_H(1+p^2)<\eta_{\max},
 \label{sharp:upper-height-splice-condition}
\end{equation}
then a common multiplier shrink adds no critical point on the
complement of the lower interval, up to the captured raw-height cap.
\end{lemma}
\begin{proof}
The bounded derivative gives \(F'>0\) on the bounded-height part
after shrinking \(m\). Put \(Z=z_*/s^2\). If
\(Z\le K_H(1+p^2)\), the displayed inequality puts every captured
height below \(\eta_{\max}/\tau\). Otherwise the intervals
\((0,\eta_{\max}/\tau]\) and \((K_H(1+p^2),Z)\) overlap
below the cap and cover it. The high-height estimate gives \(F'>0\)
on the added interval. These are restrictions of the same completed
map; scalar order retains all inputs between captured roots.
\end{proof}
